\section[Time--space complementarity]{Time--space complementarity: Weyl-pair uncertainty as a resource lower bound}
\label{sec:complementarity}

\subsection{Weyl pairs in the scan model}
In the minimal scan model on $L^2(S^1)$, let $U:=\Theta$ be the unitary shift (scan) and let $V$ be the unitary multiplication operator implementing the phase channel. For irrational slope $\alpha$ one has the Weyl commutation relation
\begin{equation}
UV=\e^{\iu\Phi}VU,
\end{equation}
with $\Phi=2\pi\alpha$.
Equivalently, $U$ and $V$ generate the irrational rotation algebra \cite{Rieffel1981}.

\subsection{A variance-type uncertainty relation}
Following the HPA formalism \cite{Ma2025HPA} (see also related Fourier-uncertainty discussions \cite{MassarSpindel2008}), define for a normalized state $|\psi\rangle$
\begin{equation}
\Delta_U^2:=1-\bigl|\langle\psi|U|\psi\rangle\bigr|^2,
\qquad
\Delta_V^2:=1-\bigl|\langle\psi|V|\psi\rangle\bigr|^2.
\end{equation}
Then $\Delta_U=0$ iff $|\psi\rangle$ is an eigenstate of $U$, and similarly for $V$.
The Weyl commutation implies a quantitative obstruction to simultaneous eigenstate-like localization.

\begin{theorem}[Weyl-pair complementarity (variance form)]
\label{thm:weyl_uncertainty}
Let $U,V$ be unitary operators satisfying $UV=\e^{\iu\Phi}VU$ and set $A=\tan(\Phi/2)$. Then for any normalized state $|\psi\rangle$,
\begin{equation}
(1+2A)\,\Delta_U^2\,\Delta_V^2 + A^2\,(\Delta_U^2+\Delta_V^2)\ge A^2.
\label{eq:weyl_variance_uncertainty}
\end{equation}
In particular, if $\Phi\not\equiv 0\ \mathrm{mod}\ 2\pi$ then $\Delta_U$ and $\Delta_V$ cannot both be arbitrarily small.
\end{theorem}

\subsection{Resource semantics: why T and S cannot both be minimized}
Equation~\eqref{eq:weyl_variance_uncertainty} becomes a resource statement once we identify $U$ with scan access and $V$ with readout phase channel.

\begin{itemize}
    \item \textbf{Scan invariance and temporal resolution.} Small $\Delta_U$ means $|\psi\rangle$ is close to an eigenstate of the scan shift, so successive scan states have large overlap. Operationally, $\Delta_U$ is the trace distance between $\rho$ and $U\rho U^{\dagger}$ for pure $\rho$ (Appendix~\ref{app:weyl_derivation}), hence small $\Delta_U$ means it is difficult to distinguish ``one tick'' from ``one tick plus a scan shift.'' In readout tasks that rely on accumulating distinguishable orbit information, such near-invariance typically forces larger scan depth.
    \item \textbf{Readout localization and phase-channel sharpness.} Small $\Delta_V$ means $|\psi\rangle$ is close to an eigenstate of the phase channel. Equivalently, the unitary conjugation $\rho\mapsto V\rho V^{\dagger}$ is hard to distinguish from $\rho$ (Appendix~\ref{app:weyl_derivation}). In the scan--readout semantics this corresponds to sharpening distinctions in the readout/phase channel, which is generally paid for by increased readout resolution (longer prefixes and/or more orthogonal cut capacity).
\end{itemize}

\noindent Appendix~\ref{app:weyl_derivation} turns this into a concrete hypothesis-testing statement: the Helstrom-optimal binary decision errors for discriminating $\rho$ from $U\rho U^{\dagger}$ and $\rho$ from $V\rho V^{\dagger}$ obey the Weyl tradeoff (Proposition~\ref{prop:weyl_binary_tradeoff}).

\paragraph{A minimal $T(\epsilon)$ lower bound (task-level).}
As a canonical example, consider the decision problem of distinguishing whether a single scan shift has occurred, i.e.\ discriminating $\rho$ from $U\rho U^{\dagger}$ for $U=\Theta$.
For pure states, Appendix~\ref{app:weyl_derivation} gives an explicit $n$-copy Helstrom formula and an inversion: to achieve error $\le\epsilon$ one needs at least
\begin{equation}
n\ \ge\ \frac{\log\!\bigl(1/(4\epsilon(1-\epsilon))\bigr)}{-\log(1-\Delta_U^2)}
\ \gtrsim\ \frac{\log(1/\epsilon)}{\Delta_U^2},
\end{equation}
where $\Delta_U^2=1-|\langle\psi|U|\psi\rangle|^2$ (Corollary~\ref{cor:copies_for_unitary_kick}).
Interpreting each repetition as one unit of scan access, this gives a concrete sense in which ``near invariance under the scan'' forces increased time cost for reliable discrimination.
\begin{remark}[Resource accounting convention]
The bound above is stated for $n$ independent repetitions (or $n$ identical copies) of the discrimination task.
In this paper we use it as a minimal operational proxy for time cost: one repetition corresponds to one unit of experimental access to the scan--readout interface.
This convention is standard in complexity-style resource accounting and is logically independent of any claims about obtaining multiple copies of a global universe state.
\end{remark}

\noindent The theorem says these objectives conflict when the scan is genuinely irrational (nontrivial $\Phi$). Thus, in HPA--$\Omega$, time--space tradeoffs are not engineering artifacts but formal consequences of (i) unitary access order and (ii) orthogonal projection readout.

This complements the reversible-computation tradeoff viewpoint (Section~\ref{sec:time}): unitarity preserves information globally, while readout discards information locally. The Weyl-pair constraint quantifies the irreducible incompatibility between ``making access cheap'' and ``making distinctions sharp''.
